\documentclass[11pt]{article}

% -------------------------------------------------
% Packages
% -------------------------------------------------
\usepackage[utf8]{inputenc}
\usepackage[T1]{fontenc}
\usepackage{amsmath, amssymb, mathtools}
\usepackage{geometry}
\usepackage{booktabs}
\usepackage{hyperref}
\usepackage{enumitem}

\geometry{margin=1in}

% -------------------------------------------------
% Document
% -------------------------------------------------
\begin{document}

\begin{center}
    {\LARGE \textbf{Exercise 2.3 — Markov Consumption Processes and Bayesian Model Comparison}}\\[0.5em]
    {\large Recursive Macroeconomic Theory}\\[0.25em]
    {\normalsize Lars Ljungqvist and Thomas J. Sargent}
\end{center}

\vspace{1em}

\section*{Problem}

Consumption is governed by an \(n\)-state Markov chain (\(P,\pi_0\)) where P is a stochastic matrix and \(\pi_0\) is an initial probability distribution. Consumption takes one of the values in the \(n\times1\) vector \(\bar{c}\). A consumer ranks stochastic processes of consumption \(t=0,1,\dots\) according to
\[
E\sum_{t=0}^\infty \beta^tu(c_t)
\]
where E is the mathematical expectation and \(u(c_t)=\frac{c^{1-\gamma}}{1-\gamma}\) for some parameter \(\gamma\geq1.\) Let \(u_i=u(\bar{c}_i)\), \(v_i=E[\sum_{t=0}^\infty\beta^tu(c_t)|x_0=\bar{e}_i]\), and \(V=Ev\), where \(\beta\in(0,1)\) is a discount factor. 

\begin{enumerate}
    \item Let u and v be the \(n\times1\) vectors whose \(i\)th component are \(u_i\) and \(v_i\), respectively. Verify the following formulas for v and V: \(v = (I-\beta P)^{-1}u\), and \(V = \sum_i\pi_{0,i}v_i\).

    \item Consider the following two Markov processes:

    Process 1: \(\pi_0=\begin{pmatrix}
        0.5 \\
        0.5
    \end{pmatrix}\), \(P=\begin{pmatrix}
        1 & 0\\
        0 & 1
    \end{pmatrix}\).

    Process 2: \(\pi_0=\begin{pmatrix}
        0.5 \\
        0.5
    \end{pmatrix}\), \(P=\begin{pmatrix}
        0.5 & .5\\
        0.5 & 0.5
    \end{pmatrix}\).

    For both Markov processes, \(\bar{c}=\begin{pmatrix}
        1 \\
        5
    \end{pmatrix}\).

    Assume that \(\gamma=2.5\), \(\beta=0.95\). Compute the unconditional discounted expected utility V for each of these processes. Which of the two processes does the consumer prefer? Redo the calculations for \(\gamma=4\). Now which process does the consumer prefer?
    \item An econometrician observes a sample of 10 observations of consumption rates for our consumer. He knows that one of the two preceding Markov processes generates the data, but he does not know which one. He assigns equal “prior probability” to the two chains. Suppose that the 10 successive observations on consumption are as follows: 1,1,1,1,1,1,1,1,1,1. Compute the likelihood of this sample under process 1 and under process 2. Denote the likelihood function \(\text{Prob}(\text{data}|\text{Model}_i)\), \(i=1,2\).

    \item Suppose that the econometrician uses Bayes’ law to revise his initial probability estimates for the two models, where in this context Bayes’ law states:

    \[
    \text{Prob}(M_i|\text{data})=\frac{\text{Prob}(\text{data}|M_i)\text{Prob}(M_i)}{\sum_j\text{Prob}(\text{data}|M_j)\text{Prob}(M_j)},
    \]
    where \(M_i\) denotes model \(i\) and the denominator is the unconditional probability of the data. After observing the data sample, what probabilities does the econometrician place on the two possible models?

    \item Repeat the calculation in part d, but now assume that the data sample is 1,5,5,1,5,5,1,5,1,5.

\end{enumerate}
\section*{Derivation of the Formulas for v and V}

We want to prove that \(v = (I-\beta P)^{-1}u\) and \(V = \sum_{i=1}^n\pi_{0i}v_i\). Recall that the Markov chain has states \(e_1, \dots, e_n\), transition matrix \(P\), and initial distribution \(\pi_0\). By definition of transition matrix, we have

\[
P_{ij}=\text{Pr}(x_{t+1}=e_j|x_t=e_i),
\]
so \(P^t\) contains the t-step transition probabilities.
Let \(u=\begin{pmatrix}
    u_1 \\
    \vdots \\
    u_n
\end{pmatrix}\), and \(u_i=u(\bar{c}_i)\).

By definition,
\[
v_i = E[\sum_{t=0}^\infty\beta^tu(c_t)|x_0=\bar{e}_i].
\]
We now derive a vector expression for \(v\). In particular, since \(\beta^t\) is deterministic,
\[
v_i = \sum_{t=0}^\infty\beta^tE[u(c_t)|x_0=\bar{e}_i],
\]
which follows from the linearity of the expectation operator. Therefore, to compute \(v_i\), it is enough to compute the conditional expected utility at every future date \(t\). Consumption takes the value \(\bar{c}_j\) whenever the Markov state is \(e_j\). Hence,
\[
u(c_t)=u_j \quad \text{if} \quad x_t=e_j.
\]
Thus \(u(c_t)\) is a function of the state. For this matter, we have a forecasting formula for a function of a Markov state. In particular,

\[
E[u(c_t)|x_0=e_i]= (P^tu)_i.
\]
We can substitute it into the definition of \(v_i\) to obtain
\[
v
=
\sum_{t=0}^{\infty}\beta^tP^tu
=
\left(I+\beta P+\beta^2P^2+\dotsb\right)u.
\]
Notice that we are dealing with a geometric series. To determine whether it converges, we have to verify that the spectral radius of \(\beta P\) is strictly less than one. Of course this is the case because P is a stochastic matrix, with its largest eigenvalue in absolute value being 1, while \(0<\beta<1\). Hence, the eigenvalues of \(\beta P\) have modulus strictly smaller than 1, so the series converges, meaning we can safely apply the matrix geometric-series formula as follows:

\[
v=(I-\beta P)^{-1}u.
\]
We now derive the formula for unconditional utility \(V\). By definition,
\[
V = Ev.
\]
More precisely, before observing \(x_0\), the initial state is random. We know that
\[
\text{Pr}(x_0=e_i)=\pi_{0i}.
\]
Conditional on \(x_0=e_i\), lifetime expected utility equals \(v_i\). Therefore, by simply applying the law of iterated expectations, we can write
\[
V=\sum_{i=1}^n\text{Pr}(x_0=e_i)E\bigg[\sum_{t=0}^\infty \beta^tu(c_t)|x_0=e_i \bigg]= \sum_{i=1}^n\pi_{0i}v_i,
\]
where the last equality followed from the definitions of \(\pi_{0i}\) and \(v_i\).

\section*{Computational Solution}

The remaining parts of the exercise are solved numerically in Python. The
implementation constructs the two transition matrices, evaluates the discounted
utility associated with each Markov process, computes the likelihood of the
observed consumption histories, and finally applies Bayes' rule to update the
probabilities assigned to the two competing models.

\subsection*{Discounted Expected Utility}

We first consider the two transition matrices

\[
P_1=
\begin{pmatrix}
1 & 0\\
0 & 1
\end{pmatrix},
\qquad
P_2=
\begin{pmatrix}
0.5 & 0.5\\
0.5 & 0.5
\end{pmatrix},
\]
with common initial distribution

\[
\pi_0=
\begin{pmatrix}
0.5\\
0.5
\end{pmatrix},
\]
and consumption vector

\[
\bar c=
\begin{pmatrix}
1\\
5
\end{pmatrix}.
\]
The Python implementation evaluates the CRRA utility vector

\[
u_i=\frac{\bar c_i^{1-\gamma}}{1-\gamma},
\]
and then computes the conditional lifetime-value vector from

\[
v=(I-\beta P)^{-1}u.
\]
The unconditional discounted expected utility is obtained as

\[
V=\pi_0'v.
\]
For $\beta=0.95$ and $\gamma=2.5$, the numerical output is

\[
u=
\begin{pmatrix}
-0.666667\\
-0.059628
\end{pmatrix}.
\]
For Process 1, Python returns

\[
v^{(1)}=
\begin{pmatrix}
-13.333333\\
-1.192570
\end{pmatrix},
\qquad
V_1=-7.262951.
\]
For Process 2, the corresponding output is

\[
v^{(2)}=
\begin{pmatrix}
-7.566471\\
-6.959432
\end{pmatrix},
\qquad
V_2=-7.262951.
\]
Thus,

\[
\boxed{V_1=V_2=-7.262951},
\]
meaning that the consumer is indifferent between the two processes from an ex ante perspective.

It is nevertheless worth noticing that the conditional value vectors are
substantially different; indeed, under Process 1, the transition matrix is the identity
matrix, so the initial consumption state is permanent. A consumer who starts
with consumption equal to $1$ remains in the low-consumption state forever,
while a consumer who starts with consumption equal to $5$ remains permanently
in the high-consumption state. This explains the large difference between the
two components of $v^{(1)}$.

Under Process 2, instead, next period's state is drawn with probability $0.5$
from either state regardless of the current state. Consequently, the initial
condition has a much smaller effect on lifetime utility, and the two components
of $v^{(2)}$ are much closer to one another.

Repeating the computation for $\gamma=4$ gives

\[
u=
\begin{pmatrix}
-0.333333\\
-0.002667
\end{pmatrix}.
\]
For Process 1,

\[
v^{(1)}=
\begin{pmatrix}
-6.666667\\
-0.053333
\end{pmatrix},
\qquad
V_1=-3.360000,
\]
whereas for Process 2,

\[
v^{(2)}=
\begin{pmatrix}
-3.525333\\
-3.194667
\end{pmatrix},
\qquad
V_2=-3.360000.
\]
Hence,

\[
\boxed{V_1=V_2=-3.360000}.
\]
Again, the consumer remains indifferent between the two processes.

The equality of the unconditional values is not accidental. Under the initial
distribution used in the exercise, both processes imply

\[
\Pr(c_t=1)=\Pr(c_t=5)=0.5
\]
at every date $t$. Therefore, the two models generate the same unconditional
distribution of consumption in every period, even though they imply radically
different serial dependence. Since preferences are additively separable over
time,

\[
E\left[\sum_{t=0}^{\infty}\beta^t u(c_t)\right]
=
\sum_{t=0}^{\infty}\beta^t E[u(c_t)],
\]
and consequently the different degree of persistence does not affect the ex ante
ranking in this particular exercise.

\subsection*{Likelihood of the First Consumption Sample}

The first observed sample is

\[
1,1,1,1,1,1,1,1,1,1.
\]
The observable consumption values are first mapped into the corresponding
zero-indexed Markov states. Since consumption equal to $1$ corresponds to state
$0$, the observed state path is
\[
(0,0,0,0,0,0,0,0,0,0).
\]
The same log-likelihood function introduced in Exercise 2.1 is then applied to
both candidate transition matrices. The Python output is

\[
\log L_1=-0.693147,
\qquad
L_1=0.500000,
\]
for Process 1, and

\[
\log L_2=-6.931472,
\qquad
L_2=0.0009765625,
\]
for Process 2.

Therefore,

\[
\boxed{
\Pr(\text{data}\mid M_1)=0.5
}
\]
and

\[
\boxed{
\Pr(\text{data}\mid M_2)=0.0009765625
}.
\]

The sample is much more likely under Process 1. This result follows directly
from the different persistence properties of the two models. Under Process 1,
once the chain starts in the low-consumption state, it remains there with
probability one. Under Process 2, instead, remaining in the same state for nine
successive transitions requires the realization of nine independent
probability-$0.5$ events.

The likelihood comparison therefore shows how observations on the time-series
pattern of consumption can distinguish between stochastic processes that have
the same unconditional distribution.

\subsection*{Bayesian Updating}

The econometrician initially assigns equal prior probabilities to the two models,

\[
\Pr(M_1)=\Pr(M_2)=0.5.
\]
Using the likelihoods computed above, the Python implementation applies Bayes'
rule and returns

\[
\Pr(M_1\mid \text{data})=0.998051,
\]
and

\[
\Pr(M_2\mid \text{data})=0.001949.
\]
Thus,

\[
\boxed{
\Pr(M_1\mid \text{data})\approx 99.8051\%
}
\]
and

\[
\boxed{
\Pr(M_2\mid \text{data})\approx 0.1949\%
}.
\]
Ten consecutive observations of low consumption therefore provide very strong
evidence in favor of the persistent model. Notice, however, that the posterior
probability of Process 2 does not become exactly zero. A sequence of ten low
observations is unlikely under Process 2, but it remains possible.

This illustrates an important role of likelihood-based inference; that is, the
econometrician does not compare the models according only to whether a sample
is possible under each one, but according to how much probability each model
assigns to the observed history.

\subsection*{Second Consumption Sample}

Finally, we consider the second observed history,

\[
1,5,5,1,5,5,1,5,1,5,
\]
which corresponds to the state path

\[
(0,1,1,0,1,1,0,1,0,1).
\]
The Python output gives

\[
L_1=0
\]
for Process 1, while

\[
L_2=0.0009765625
\]
for Process 2.

Consequently, Bayes' rule yields

\[
\boxed{
\Pr(M_1\mid\text{data})=0
}
\]
and

\[
\boxed{
\Pr(M_2\mid\text{data})=1
}.
\]
The interpretation is immediate. Process 1 has transition matrix

\[
P_1=
\begin{pmatrix}
1&0\\
0&1
\end{pmatrix},
\]
so transitions between states have probability zero. The first observed
transition from consumption $1$ to consumption $5$ is therefore sufficient to
rule out Process 1 completely.

This example highlights the distinction between an observation that is merely
unlikely and one that is impossible under a model. The first data sample strongly
favored Process 1 but did not rule out Process 2. The second sample, by contrast,
contains transitions to which Process 1 assigns exactly zero probability, so its
posterior probability collapses to zero.

\section*{Computational Extensions}

The exercise can be extended in several directions that provide additional
economic interpretation. In particular, the two Markov processes have identical
one-period unconditional consumption distributions but very different dynamic
properties. This makes them useful for studying the distinction between
cross-sectional risk and persistence over time.

\subsection*{Sequential Bayesian Learning}

The previous Bayesian calculation uses the complete sample at once. A natural
extension is to update the posterior probabilities sequentially as each new
observation becomes available.

For the sample

\[
1,1,1,1,1,1,1,1,1,1,
\]
the Python implementation produces the following posterior probabilities for
Process 1:

\[
\begin{array}{c|cc}
\toprule
\text{Number of observations} &
\Pr(M_1\mid\text{data}) &
\Pr(M_2\mid\text{data})\\
\midrule
1  & 0.500000 & 0.500000\\
2  & 0.666667 & 0.333333\\
3  & 0.800000 & 0.200000\\
4  & 0.888889 & 0.111111\\
5  & 0.941176 & 0.058824\\
10 & 0.998051 & 0.001949\\
\bottomrule
\end{array}
\]
The first observation does not distinguish between the models because both
assign probability $0.5$ to starting in the low-consumption state. However,
every additional low-to-low transition provides new evidence in favor of the
persistent process.

This sequential interpretation is particularly relevant in macroeconomic
applications. Suppose, for example, that the low-consumption state represents a
household experiencing a negative income shock. A single low observation might
be interpreted as a temporary disturbance. A long sequence of low observations,
instead, progressively raises the probability that the household is affected by
a highly persistent or permanent shock.

The same logic applies to real-time macroeconomic inference, where policymakers
continuously update beliefs about whether an observed slowdown is temporary or
persistent as new data arrive.

\subsection*{Persistence and Autocorrelation}

A second extension computes the theoretical autocorrelation function of
consumption.

For Process 1, Python returns

\[
\rho_k=1
\qquad
\text{for }k=1,\ldots,5.
\]
For Process 2, the output is

\[
\rho_k=0
\qquad
\text{for }k=1,\ldots,5.
\]
Thus,

\[
\boxed{
\operatorname{Corr}_{P_1}(c_t,c_{t+k})=1
}
\]
while

\[
\boxed{
\operatorname{Corr}_{P_2}(c_t,c_{t+k})=0
}
\]
for every positive lag.

The difference follows directly from the transition matrices. Under Process 1,
the state never changes, so knowing current consumption perfectly predicts all
future consumption. Under Process 2, both rows of the transition matrix are
identical, so the distribution of next period's consumption does not depend on
the current state.

This distinction can be interpreted in terms of permanent versus transitory
income risk. Two households may face exactly the same probability of having low
income in any given year while facing very different risks over their lifetimes.
If low income is highly persistent, current economic conditions contain much
more information about future conditions than when shocks are purely
transitory.

Such persistence plays an important role in models of precautionary saving,
consumption smoothing, incomplete markets, and household insurance.

\subsection*{Expected Duration of Consumption States}

Another useful measure of persistence is the expected duration of a spell in a
given state. For a Markov state with persistence probability $P_{ii}$, the
expected duration is

\[
E[D_i]=\frac{1}{1-P_{ii}},
\]
provided that $P_{ii}<1$.

For Process 1, Python returns

\[
E[D_1]=E[D_2]=\infty.
\]
Both states are absorbing, so once a consumer enters either state, the consumer
remains there forever.

For Process 2,

\[
P_{11}=P_{22}=0.5,
\]
and the numerical output is therefore

\[
\boxed{
E[D_1]=E[D_2]=2
}
\]
periods.

The duration statistic provides an intuitive interpretation of the two processes.
If the low-consumption state is associated with unemployment or a negative
income shock, Process 1 represents a permanent shock, whereas Process 2 implies
an expected low-income spell of only two periods.

The unconditional probability of low consumption is identical in the two
models, but the welfare and policy implications of a temporary and a permanent
income shock can be very different once households are allowed to save,
borrow, or insure themselves.

\subsection*{Certainty-Equivalent Consumption}

The numerical values of lifetime utility are difficult to interpret directly.
A useful welfare transformation is therefore the certainty-equivalent level of
consumption. Let $c^{CE}$ denote the constant level of consumption that generates
the same discounted lifetime utility $V$:

\[
V=
\sum_{t=0}^{\infty}
\beta^t u(c^{CE}).
\]
The Python computation gives, for $\gamma=2.5$,

\[
\boxed{
c^{CE}=1.499283
}
\]
for both processes.

For $\gamma=4$, the result is

\[
\boxed{
c^{CE}=1.256579
}
\]
for both processes.

As expected, the certainty equivalent falls as $\gamma$ rises. Greater curvature
of the utility function makes the consumer more sensitive to low-consumption
realizations, lowering the constant amount of consumption that the consumer
regards as equivalent to the risky process.

The equality of the certainty equivalents across the two models mirrors the
equality of their ex ante expected utilities. Nevertheless, this result should
not be interpreted as implying that persistence is generally irrelevant for
welfare. It is a consequence of the particular preference specification used in
the exercise, together with the fact that both processes have the same
unconditional consumption distribution at every date.

\section*{Discussion}

The numerical results reveal a useful distinction between the marginal
distribution of a stochastic process and its dynamic dependence structure.

Both Markov processes satisfy

\[
\Pr(c_t=1)=\Pr(c_t=5)=0.5
\]
at every date. Consequently, they generate the same expected period utility and,
under time-additive expected utility, the same ex ante discounted lifetime
utility.

However, their time-series properties are completely different. Under Process 1,
consumption is perfectly persistent and each state is absorbing. Under Process 2,
consumption is serially independent and expected state durations are finite.

The likelihood and Bayesian calculations demonstrate that these differences can
be identified from observed histories even though they cannot be detected from
the unconditional distribution alone. A long spell in one state provides evidence
in favor of a persistent process, while the observation of even a single switch
immediately rules out a model in which transitions are impossible.

The exercise therefore illustrates two complementary uses of Markov chains.
From the consumer's perspective, the transition matrix determines forecasts and
conditional lifetime values. From the econometrician's perspective, the same
transition matrix determines the probability of observed histories and therefore
provides the basis for likelihood-based model comparison and Bayesian learning.


\section*{Reference}

Ljungqvist, L. and Sargent, T. J. (2018).
\textit{Recursive Macroeconomic Theory},
4th edition, MIT Press.

\end{document}
